\subsection{矩阵列点积：Cache 策略}

\texttt{matrix\_dot\_cache} 与 Naive 在数学上等价，但重排了访存局部性（memory locality）友好的执行顺序。
在实现上，外层遍历行 \(i\)，并把列方向按块宽 \(B\) 分块（tiling）：
\begin{equation}
 o_j \leftarrow o_j + A_{ij}v_i,\quad
j\in[b,\min(b+B,n)).
\end{equation}

该重排不改变计算结果：
\begin{equation}
\forall j,\quad o_j=\sum_{i=1}^{m}A_{ij}v_i,
\end{equation}
但通常可改善缓存命中（cache hit）与预取行为（prefetch behavior）。

复杂度仍为
\begin{equation}
T(m,n,B)=\Theta(mn),
\end{equation}
差异主要体现在常数项（constant factor）。
当 \(B\) 与缓存层级匹配时，实测性能往往优于 Naive；
当 \(B\) 过小或过大时，分块调度开销与缓存压力可能抵消收益。

参数合法性条件为
\begin{equation}
B>0,
\end{equation}
否则抛出 \texttt{std::invalid\_argument}。

从策略设计角度，Cache 版本体现了“先保持语义不变，再优化数据路径（data path）”的原则，
是本项目中最典型的“零语义风险优化”（semantics-preserving optimization）。
